\documentclass[conference]{IEEEtran}
\usepackage{cite}
\usepackage{amsmath,amssymb}
\usepackage{url}
\usepackage{balance}
\usepackage[hidelinks]{hyperref}

\newcommand{\reporttitle}{Technical Debt in Large Language Model Application Development and Deployment}
\newcommand{\reportabstract}{Large language model applications introduce maintenance obligations in prompts, generation settings, software integrations, and production infrastructure. This report critically reviews two complementary accounts of these obligations. PromptDebt provides empirical evidence from developer comments and identifies five categories of self-admitted technical debt. A deployment-oriented study synthesizes engineering challenges involving memory management, scheduling, model optimization, and system dependencies. The comparison indicates that debt must be examined across application artifacts and their operational context. It also reveals a methodological gap between identifying documented shortcomings and measuring their consequences. We argue that prompts warrant explicit maintenance practices, with evaluation effort proportional to application risk. Future research should relate recorded debt to maintenance effort, response quality, and operating cost, and assess the effectiveness of proposed mitigation practices.}
\newcommand{\reportkeywords}{large language models, technical debt, prompt engineering, software maintenance}

\title{\reporttitle}
\author{\IEEEauthorblockN{[Your name] and [Teammate names]}
\IEEEauthorblockA{SOEN 6431\\Student IDs: [Your ID] and [Teammate IDs]}}

\begin{document}
\onecolumn
\thispagestyle{empty}
\begin{center}
\vspace*{0.7in}
{\LARGE\bfseries \reporttitle\par}
\vspace{0.35in}
{\large SOEN 6431\\Literature review report\par}
\vspace{0.3in}
Chen Qian \quad 27867808\\[0.08in]
[Teammate names] \quad [Teammate student IDs]
\end{center}

\vspace{0.3in}
\begin{center}
\textbf{Contributions}
\end{center}
\noindent\begin{tabular}{|p{0.16\textwidth}|p{0.19\textwidth}|p{0.56\textwidth}|}
\hline
\textbf{Student ID} & \textbf{Name} & \textbf{Contribution} \\
\hline
27867808 & Chen Qian & Reviewed PromptDebt and Navigating Challenges and Technical Debt in Large Language Models Deployment; drafted Sections~I--V and the references for group revision. \\
\hline
[Teammate ID] & [Teammate name] & [Reviewed paper and contributed sections to be added.] \\
\hline
\end{tabular}

\vspace{0.3in}
\noindent\textbf{Abstract}\par
\smallskip
\noindent\reportabstract

\medskip
\noindent\textbf{Keywords:} \reportkeywords

\clearpage
\twocolumn
\setcounter{page}{2}
\pagestyle{plain}
\maketitle
\thispagestyle{plain}

\begin{abstract}
\reportabstract
\end{abstract}
\begin{IEEEkeywords}
\reportkeywords
\end{IEEEkeywords}

% GROUP REVISION: Integrate teammates' research questions and findings into this motivation.
\section{Introduction}\label{sec:introduction}
Large language models (LLMs) support applications that generate text, answer questions, and process information expressed in natural language. Their integration also creates maintenance responsibilities involving prompts, generation settings, software frameworks, and production infrastructure. As requirements and deployment conditions change, decisions made to deliver initial functionality can increase the effort required to sustain the application. The reviewed literature examines these obligations through the concept of technical debt~\cite{aljohani2025promptdebt,menshawy2024deployment}.

Technical debt concerns the future consequences of development and design decisions that leave deficiencies unresolved or defer necessary improvements. In LLM applications, these decisions may involve incomplete prompt templates, hardcoded configuration, fragile framework dependencies, or inadequate provision for production workloads. Their significance lies in the resulting maintenance burden. Identifying debt therefore requires attention to the decision and its consequences, alongside the technical difficulty of the underlying task~\cite{aljohani2025promptdebt,menshawy2024deployment}.

This report examines technical debt in LLM application development and deployment through a focused critical review of two publications. Aljohani and Do's PromptDebt study investigates self-admitted technical debt (SATD), which developers acknowledge in source-code comments. It provides empirical evidence of deficiencies in prompts, generation settings, framework integrations, cost management, and learning strategies~\cite{aljohani2025promptdebt}. Menshawy, Nawaz, and Fahmy examine the surrounding production system, including memory management, scheduling, model optimization, integration, and hidden dependencies~\cite{menshawy2024deployment}. These perspectives connect application-level decisions with the operational conditions under which their consequences emerge.

The papers also differ in their evidential basis. PromptDebt derives categories and occurrence counts from a selected repository corpus, whereas the deployment paper synthesizes published research and engineering experience. Their findings must therefore be interpreted according to what each method can establish. Recorded SATD indicates acknowledged deficiencies within the analyzed artifacts. A narrative account of deployment challenges identifies mechanisms and candidate responses that require subsequent evaluation. Comparing the papers makes the relationship between documented shortcomings and operational consequences a central research question~\cite{aljohani2025promptdebt,menshawy2024deployment}.

The review contributes an account of the maintenance concerns identified in the two papers and an assessment of the evidence supporting their conclusions. It also examines whether prompts should receive the same maintenance discipline as conventional source code. This question connects the empirical observations about prompt deficiencies with the broader dependencies of deployed systems. The discussion considers both the value of explicit maintenance practices and the costs of applying them during experimentation. It then proposes research directions that could connect debt identification with measurable outcomes.

The remainder of the report is organized as follows. Section~\ref{sec:background} introduces the relevant concepts. Section~\ref{sec:reviews} reviews each paper's motivation, methodology, findings, and conclusions. Section~\ref{sec:discussion} evaluates the controversial statement and identifies future research directions. Section~\ref{sec:conclusions} presents the conclusions.

% GROUP REVISION: Extend these definitions only where teammates' papers require additional concepts.
\section{Background}\label{sec:background}
\subsection{LLM applications and prompt engineering}
An LLM application combines a language model with software that constructs inputs, manages interactions, and processes outputs. The model may be accessed through a hosted application programming interface (API) or served on infrastructure managed by the application provider. These arrangements allocate maintenance responsibilities differently. API integration emphasizes prompts, configuration, and interfaces, while self-managed serving introduces direct responsibility for memory, computation, and request scheduling~\cite{aljohani2025promptdebt,menshawy2024deployment}.

Prompts specify tasks through instructions, contextual information, examples, and expected output formats. Instruction prompts describe the requested behavior, few-shot prompts provide examples, and documentation prompts supply source material. Prompt engineering concerns the construction and revision of these inputs. Generation settings, including temperature and output-token limits, provide additional controls. Orchestration frameworks can organize prompt templates and process model responses, creating dependencies between natural-language instructions, configuration, and conventional application code~\cite{aljohani2025promptdebt}.

\subsection{Technical debt and developer acknowledgments}
Technical debt describes a continuing maintenance obligation associated with a development or design decision. SATD is the subset explicitly acknowledged by developers, commonly through comments indicating incomplete work or a necessary improvement. Such comments offer observable evidence of recognized deficiencies, but their absence does not establish the absence of debt. Interpretation depends on the artifacts collected and the procedure used to identify relevant acknowledgments~\cite{aljohani2025promptdebt}.

PromptDebt distinguishes debt related to prompts, generation settings, integrated frameworks, costs, and learning strategies. These categories describe how developers use and adapt LLM capabilities. The deployment perspective also considers dependencies among models, data, infrastructure, and users. Entanglement refers to interactions through which changes to one component affect behavior elsewhere. Hidden feedback loops occur when generated outputs influence subsequent inputs or data used for later training. Both mechanisms can complicate maintenance across component boundaries~\cite{aljohani2025promptdebt,menshawy2024deployment}.

\subsection{Adaptation and production serving}
Fine-tuning adapts model parameters through additional training. Retrieval-augmented generation (RAG) supplies retrieved information as context during inference. RAG can provide access to a changing knowledge source while introducing obligations to maintain the retrieval corpus and its embeddings. Prompt design and evaluation remain necessary because the usefulness of retrieved context depends on its relevance and subsequent use in the response~\cite[Sec.~2]{menshawy2024deployment}.

Serving also requires management of the model parameters and key-value (KV) cache, which retains attention information during generation. Batching groups requests to improve hardware utilization, while continuous batching permits the active group to change as requests progress. Quantization reduces numerical precision, pruning removes selected parameters, and parallelism distributes computation across devices. These approaches address different resource constraints and involve trade-offs among memory use, response quality, latency, and implementation complexity. Their maintenance implications therefore depend on the model, workload, and surrounding system~\cite[Secs.~2.1 and 4]{menshawy2024deployment}.

% INDIVIDUAL CONTRIBUTION: Both reviews share the agreed maximum of two IEEE-formatted pages.
% GROUP REVISION: Insert teammates' individual reviews after these, retaining the four required headings.
\section{Individual Paper Reviews}\label{sec:reviews}
\subsection{PromptDebt: A Comprehensive Study of Technical Debt Across LLM Projects}
\subsubsection{Motivation}
Aljohani and Do~\cite{aljohani2025promptdebt} investigate how application development through LLM APIs gives rise to SATD. Earlier research concentrated on traditional software and machine learning pipelines, leaving prompt construction, generation settings, orchestration, and token-based costs less fully characterized. The study addresses this gap by examining the distribution of recorded debt, identifying LLM-specific categories, and analyzing the prompting techniques represented in debt-associated examples. Its contribution is an empirical account of maintenance concerns expressed by developers.

\subsubsection{Methodology}
The study combines automated detection with manual categorization. PromptSet, compiled in January 2024 from 37,944 GitHub repositories, supplies an initial corpus of 93,142 Python files. Provider and framework filtering reduces this collection to 91,377 files, of which 70,884 contain single-line comments. Cleaning reduces 1,302,419 extracted comments to 340,840, and SATDDetector with complementary keywords identifies 6,042 SATD instances~\cite[Secs.~3.1--3.2]{aljohani2025promptdebt}.

The authors manually categorize a stratified sample of 998 comments. Two investigators label independently, obtain Cohen's kappa of 0.74, and resolve disagreements with a third investigator. A separate prompt-analysis stage describes 155 retained prompts with explicit content. An LLM classifies their techniques, and manual inspection judges 37 of 40 sampled classifications correct. This validation concerns the inspected prompt classifications rather than the complete SATD detection procedure~\cite[Secs.~3.3--3.5]{aljohani2025promptdebt}.

\subsubsection{Findings}
Detected SATD constitutes approximately 1.77\% of the cleaned comments and occurs in 3,630 of the 70,884 files containing comments~\cite[Sec.~3.2]{aljohani2025promptdebt}. Within the 998-comment manual sample, the five LLM-specific categories comprise prompt debt with 66 instances, hyperparameter debt with 45, integrated-framework debt with 43, cost debt with 21, and learning debt with 6. Prompt debt is the most frequent of these categories, accounting for 6.61\% of the full manual sample~\cite[Table~3]{aljohani2025promptdebt}.

The examples associate prompt debt with incomplete templates, hardcoded context, unclear instructions, excessive length, and deferred output-format specifications. Framework debt includes postponed improvements to prompt construction and response parsing. The analysis also relates prompt length and output limits to token consumption, suggesting interactions between prompt, configuration, and cost concerns~\cite[Secs.~4.2 and 5]{aljohani2025promptdebt}.

The OpenAI-only category contains 3,292 of the 6,042 SATD comments, while OpenAI with LangChain contains 1,232~\cite[Table~4]{aljohani2025promptdebt}. These counts describe a corpus with unequal provider representation and comment volumes. The prompt analysis reports instruction, few-shot, and documentation techniques most frequently among the debt-associated examples, with concerns involving clarity, example relevance, and document length~\cite[Table~5]{aljohani2025promptdebt}. This distribution cannot establish comparative susceptibility without the corresponding prevalence of each technique in the wider prompt population.

\subsubsection{Conclusions}
The authors recommend reusable templates, task-specific evaluation of generation settings, structured response handling, and deliberate model selection~\cite[Secs.~5--7]{aljohani2025promptdebt}. The study supports treating prompts and configuration as maintained software artifacts. Its scope nevertheless remains restricted to acknowledged debt in selected Python repositories, and the use of a detector developed for traditional software introduces classification uncertainty. Comment counts also provide no direct measure of severity or repayment effort.

Reporting inconsistencies further limit the precision of the prompt analysis. Table~3 percentages differ from the accompanying prose, while prompt-technique percentages and sample totals vary across Sections~3.5 and 4.3 and Table~5~\cite{aljohani2025promptdebt}. Accordingly, this review uses the Table~3 counts and denominator and treats the prompt-technique ranking qualitatively. The taxonomy provides a basis for further investigation, while the effectiveness of its proposed mitigation practices remains to be evaluated.

\subsection{Navigating Challenges and Technical Debt in Large Language Models Deployment}
\subsubsection{Motivation}
Menshawy, Nawaz, and Fahmy~\cite{menshawy2024deployment} examine the engineering responsibilities surrounding production LLM deployment. Model capabilities depend on infrastructure, interfaces, data, and evaluation processes that must remain effective as workloads and requirements change. Building on the concept of hidden technical debt in machine learning, the paper considers how deployment decisions create dependencies and continuing maintenance obligations. Its purpose is to connect these concerns with possible engineering responses, extending the analysis beyond individual application artifacts.

\subsubsection{Methodology}
The paper presents a narrative synthesis of published research and engineering experience. It organizes challenges into technical, operational, interaction and system-design, and ethical dimensions. It then examines entanglement and hidden feedback loops before discussing approaches to serving optimization. The acknowledgments relate the discussion to the authors' research and AI engineering practice~\cite[Secs.~2--4 and acknowledgments]{menshawy2024deployment}.

The approach integrates concerns that are often considered separately, but the publication reports no systematic search and selection protocol or original comparative experiment. Its figures illustrate concepts and draw on cited research. The resulting evidence supports a conceptual account of deployment problems and candidate responses. It does not provide estimates of their prevalence, severity, or comparative remediation costs.

\subsubsection{Findings}
Memory management forms a central part of the analysis. Model parameters and the KV cache consume GPU memory, and inefficient cache allocation can restrict concurrency. The paper discusses cache sharing, pruning, quantization, and semantic caching as possible responses. Quantization reduces representation precision, while semantic caching reuses responses for sufficiently similar queries. Their suitability depends on the application's resource and response-quality requirements~\cite[Secs.~2.1 and 4.1--4.2]{menshawy2024deployment}.

Scheduling and parallel execution address additional serving constraints. Continuous batching allows new requests to enter an active batch as generation progresses. Distributing computation across devices can support larger models, while introducing communication and synchronization overhead. Custom GPU kernels offer further optimization through specialized implementation. These approaches constitute alternative engineering choices whose combined benefits and maintenance costs require evaluation under the intended workload~\cite[Secs.~4.3--4.5]{menshawy2024deployment}.

The operational discussion extends to integration with existing interfaces and data pipelines, evaluation of prompt changes, and maintenance of retrieval data. RAG provides access to external information but introduces dependencies on corpus quality and embedding updates. Hallucination, privacy, bias, and accountability add further requirements for evaluation and monitoring~\cite[Secs.~2.1--2.3]{menshawy2024deployment}.

The conceptual analysis centres on entanglement and hidden feedback loops. Changes to models, data, or their environment can have effects across component boundaries, while generated outputs can influence later data and reinforce errors. These mechanisms explain why maintenance obligations may be distributed across the deployed system~\cite[Sec.~3]{menshawy2024deployment}. In this review, a deployment challenge constitutes debt when a development or operational decision leaves a deficiency unresolved or defers necessary improvement.

\subsubsection{Conclusions}
The authors advocate continuous monitoring, careful prompt engineering, responsible data practices, and infrastructure suited to the deployment context~\cite[conclusion]{menshawy2024deployment}. Their principal contribution is a systems perspective that connects model serving with application and data maintenance. This perspective complements PromptDebt's observations by identifying dependencies that may remain invisible in developer comments.

The breadth of the discussion is accompanied by limited empirical validation. The paper offers neither an operational measurement procedure for debt nor a controlled assessment of its proposed remedies. Its recommendations should therefore be understood as candidates for workload-specific evaluation. A useful extension would examine whether an optimization reduces total maintenance effort or redistributes that effort to new infrastructure and dependencies.

% GROUP REVISION: Revisit both positions and the synthesis after incorporating teammates' evidence.
\section{Controversial Statement and Discussion}\label{sec:discussion}
\subsection{Statement and competing arguments}
\textit{Prompt engineering should receive the same maintenance discipline as conventional source code.} We propose this statement for discussion because prompts influence application behavior while depending on model and data conditions that may change independently. Maintenance discipline here includes ownership, version history, review, and repeatable evaluation. The controversy concerns whether equivalent discipline is appropriate across experimentation and production, and whether conventional practices adequately address these dependencies.

The supporting position follows from the recurring deficiencies documented in PromptDebt, including incomplete templates, ambiguous instructions, and deferred output-format specifications~\cite{aljohani2025promptdebt}. Treating prompts and generation settings as maintained artifacts would make changes attributable and provide a basis for systematic comparison of behavior. Review and representative evaluation cases could also support the identification of unintended changes. These are proposed implications of the evidence, rather than measured benefits of interventions evaluated in the study.

The opposing position concerns cost and proportionality. Extensive review and evaluation of every exploratory prompt may consume time and inference resources before requirements stabilize. Moreover, prompt-level checks provide incomplete coverage when behavior depends on retrieval, model updates, and serving conditions. The deployment paper's analysis of system dependencies supports this concern~\cite{menshawy2024deployment}. Our assessment favors consistent ownership and versioning, with evaluation effort scaled to application risk and the deployed context. Equivalent responsibility need not imply identical procedures for every artifact.

% GROUP REVISION: Extend these research priorities using the limitations in teammates' reviews.
\subsection{Future research directions}
The methodological gap between recorded SATD and operational consequences motivates longitudinal research. Studies could follow prompt and configuration changes alongside response quality, token consumption, incidents, and remediation effort. Separating application changes from provider and retrieval updates would help identify which dependencies contribute to subsequent maintenance work. This would extend PromptDebt's occurrence-based evidence toward measures of impact~\cite{aljohani2025promptdebt}.

A second priority is controlled evaluation of prompt maintenance practices within deployed applications. Comparing templates, reviewed changes, and regression evaluations under shared tasks could test whether they reduce later rework. The evaluation should include their own costs and interactions with retrieval and serving infrastructure, addressing the system-level concerns identified by Menshawy et al.~\cite{menshawy2024deployment}.

% GROUP REVISION: Update the conclusions to synthesize all papers once the group report is complete.
\balance
\section{Conclusions}\label{sec:conclusions}
The two papers identify complementary dimensions of technical debt in LLM applications. PromptDebt provides empirical evidence of acknowledged deficiencies in prompts, generation settings, frameworks, costs, and learning strategies. The deployment paper examines dependencies involving infrastructure, integration, data, and feedback processes~\cite{aljohani2025promptdebt,menshawy2024deployment}. Together, they support examining maintenance across application artifacts and the systems in which those artifacts operate. Their different methods also limit the conclusions available: comment frequency does not establish operational severity, and proposed deployment remedies require empirical validation.

The review supports explicit ownership and change management for prompts, coupled with evaluation proportional to application risk. Such practices should be assessed within the complete deployed system. Future research should connect documented debt with remediation effort, response quality, and operating cost, then evaluate whether particular maintenance practices improve those outcomes. This would provide a stronger empirical basis for prioritizing debt repayment in LLM applications.

\bibliographystyle{IEEEtran}
\bibliography{references}
\end{document}
